%=================================
\begin{name}
	{\tenchude}
	{\tendethi}
	{\tentruong}
	{\thoigian}
\end{name}
\Opensolutionfile{ansbook}[ans/ansbook-DE-PNL-02]
\TN
\Opensolutionfile{ans}[ans/ans-DE-PNL-02-T]
\begin{ex}%Câu 1
	Cho hàm số $y=f(x)$ có bảng biến thiên như sau:
	\begin{center}
		\begin{tikzpicture}[>=stealth]
			\tkzTabInit[espcl=2.5,lgt=1.5,nocadre=false]
			{$x$/0.7,$f'(x)$/0.7,$f(x)$/2.1}
			{$-\infty$,$-1$,$1$,$+\infty$}
			\tkzTabLine{,+,0,-,0,+,}
			\tkzTabVar{-/$-\infty$,+/$2$,-/$-2$,+/$+\infty$}
		\end{tikzpicture}
	\end{center}
	Hàm số đã cho nghịch biến trên khoảng nào dưới đây?
	\choice
	{$\left(-2;2\right)$}
	{\True $\left(-1;1\right)$}
	{$\left(-2;1\right)$}
	{$\left(-1;+\infty\right)$}
	\loigiai{
		Từ bảng biến thiên ta suy ra hàm số nghịch biến trên khoảng $\left(-1;1\right)$.}
\end{ex}

\begin{ex}%Câu 2
	Cho cấp số nhân $\left(u_n\right)$ với $u_1=6$ và $u_2=-12.$ Công bội $q$ của cấp số nhân đã cho là
	\choice
	{$q=-\dfrac{1}{2}$}
	{\True $q=-2$}
	{$q=-18$}
	{$q=-6$}
	\loigiai{
		Ta có $u_2=u_1.q\Leftrightarrow-12=6.q\Leftrightarrow q=-2$.}
\end{ex}

\begin{ex}%Câu 3
	Trong không gian $Oxyz$, cho hai điểm $A\left(1;1;-2\right)$ và $B\left(3;-1;2\right)$. Tọa độ của vectơ $\overrightarrow{BA}$ là
	\choice
	{$\left(2;-2;4\right)$}
	{$\left(2;0;0\right)$}
	{$\left(1;-1;2\right)$}
	{\True $~\left(-2;2;-4\right)$}
	\loigiai{
		Tọa độ của vectơ $\overrightarrow{BA}$ là: $\overrightarrow{BA}=\left(x_A-x_B;y_A-y_B;z_A-z_B\right)=\left(-2;2;-4\right)$ .}
\end{ex}

\begin{ex}%Câu 4
	Tính thể tích vật thể tròn xoay khi quay hình phẳng giới hạn bởi các đường cong $y=\sqrt{e^x-x},$ $y=0$, $x=1$, $x=2$ xung quanh trục $Ox$ là
	\choice
	{\True $\pi\left(e^2-e-\dfrac{3}{2}\right)$}
	{$e^2-e-\dfrac{5}{2}$}
	{$\pi\left(e^2-e-\dfrac{5}{2}\right)$}
	{$e^2-e-\dfrac{3}{2}$}
	\loigiai{
		Hàm số $y=\sqrt{e^x-x}$ liên tục và không âm trên đoạn $\left[1;2\right]$ nên thể tích vật thể tròn xoay khi quay hình phẳng giới hạn bởi các đường $y=\sqrt{e^x-x},$ $y=0,x=1,x=2$ xung quanh trục Ox là:\\
		$V=\pi\displaystyle\int\limits_1^2\left(\sqrt{e^x-x}\right)^2dx=\pi\displaystyle\int\limits_1^2\left(e^x-x\right)dx=\left.\pi\left(e^x-\dfrac{1}{2}{x^2}\right)\right|_1^2$ $=\pi\left[\left(e^2-\dfrac{1}{2}{2^2}\right)-\left(e^1-\dfrac{1}{2}{1^2}\right)\right]=\pi\left(e^2-e-\dfrac{3}{2}\right)$ .}
\end{ex}

\begin{ex}%Câu 5
	Với mọi số thực dương a, $\log_3\left(27a\right)-\log_3a$ bằng
	\choice
	{$\log_3\left(26a\right)$}
	{$9$}
	{\True 3}
	{$3-2\log_3a$}
	\loigiai{
		Ta có: $\log_3\left(27a\right)-\log_3a=\log_327+\log_3a-\log_3a=3$ .}
\end{ex}

\begin{ex}%Câu 6
	Trong không gian $Oxyz$ , cho đường thẳng $d:\dfrac{x-1}{4}=\dfrac{-y}{2}=\dfrac{z+2}{-6}$. Vectơ nào dưới đây là một vectơ chỉ phương của $d$?
	\choice
	{$\overrightarrow{u_2}=\left(2;-1;3\right)$}
	{$\overrightarrow{u_1}=\left(4;2;-6\right)$}
	{\True $\overrightarrow{u_3}=\left(-2;1;3\right)$}
	{$\overrightarrow{u_4}=\left(1;0;2\right)$}
	\loigiai{
		Phương trình đường thẳng $d:\dfrac{x-1}{4}=\dfrac{-y}{2}=\dfrac{z+2}{-6}$ được viết lại $\dfrac{x-1}{4}=\dfrac{y}{-2}=\dfrac{z+2}{-6}$\\
		Suy ra một vectơ chỉ phương của $d$ là $\overrightarrow{u_3}=\left(-2;1;3\right)$.}
\end{ex}

\begin{ex}%Câu 7
	Tiệm cận ngang của đồ thị hàm số $y=\dfrac{2x-3}{x+1}$ là đường thẳng có phương trình:
	\choice
	{$y=-1$}
	{$x=-1$}
	{\True $y=2$}
	{$x=2$}
	\loigiai{
		Tập xác định: $D=\mathbb{R}\setminus\left\{-1\right\}$ thì ta có $\underset{x\to+\infty}{\mathop{\lim y}}=\underset{x\to+\infty}{\lim}\,\dfrac{2x-3}{x+1}=2;\underset{x\to-\infty}{\mathop{\lim y}}\,=\underset{x\to-\infty}{\lim}\,\dfrac{2x-3}{x+1}=2\,$ .\\
		Suy ra $y=2$ là tiệm cận ngang của đồ thị hàm số.}
\end{ex}

\begin{ex}%Câu 8
	Trong không gian $Oxyz,$ cho mặt cầu $(S)$ có tâm $I\left(0;-2;1\right)$ và bán kính $R=5$. Phương trình của $(S)$ là
	\choice
	{\True $x^2+\left(y+2\right)^2+\left(z-1\right)^2=25$}
	{$x^2+\left(y-2\right)^2+\left(z+1\right)^2=25$}
	{$x^2+\left(y+2\right)^2+\left(z-1\right)^2=5$}
	{$x^2+\left(y-2\right)^2+\left(z+1\right)^2=5$}
	\loigiai{
		Mặt cầu tâm $I\left(a;b;c\right)$ và bán kính $R$ có phương trình là $\left(x-a\right)^2+\left(y-b\right)^2+\left(z-c\right)^2=R^2$.}
\end{ex}

\begin{ex}%Câu 9
	Công thức tính thể tích của một khối trụ có bán kính đáy là $R$ và chiều cao $h$ là
	\choice
	{$V=2\pi{R^2}h$}
	{$V=\dfrac{4}{3}\pi{R^2}h$}
	{$V=\dfrac{1}{3}\pi{R^2}h$}
	{\True $V=\pi{R^2}h$}
	\loigiai{
		Công thức tính thể tích khối trụ là $V=\pi{R^2}h$}
\end{ex}

\begin{ex}%Câu 10
	Diện tích hình phẳng gạch sọc trong hình vẽ bên dưới bằng
	\begin{center}
		\begin{tikzpicture}[>=stealth,thick]
			\tikzset{every node/.style={scale=0.9},y=.7cm}
			\draw[->] (-3.1,0)--(5.1,0) node[below left] {$x$};
			\draw[->] (0,-1.1)--(0,9.1) node[below left] {$y$};
			\draw (0,0) node [below left] {$O$};
			\foreach \x/\nx in {1/1,3/3}
			\draw[thin] (\x,1pt)--(\x,-1pt) node [below] {$\nx$};
			\foreach \y/\ny in {2/2,8/8}
			\draw[thin] (1pt,\y)--(-1pt,\y) node [left] {$\ny$};
			\draw[dashed,thin](1,0)--(1,2)--(0,2) (1,2)--(3,2);
			\draw[dashed,thin](3,0)--(3,8)--(0,8);
			\draw[pattern=vertical lines,pattern color=red] (1,2)--(3,2)--(3,8)--plot[samples=200,domain=3:1,smooth](\x,{2^(\x)});
			\draw (2,6) node{$y=2^x$};
			\begin{scope}
				\clip (-3,-1) rectangle (5,9);
				\draw[samples=200,domain=-3:3.1,smooth] plot (\x,{2^(\x)});
			\end{scope}
		\end{tikzpicture}
	\end{center}
	\choice
	{\True $\displaystyle\int\limits_1^3\left(2^x-2\right)\,\mathrm{\,d}x$}
	{$\displaystyle\int\limits_1^3\left(2^x+2\right)\,\mathrm{\,d}x$}
	{$\displaystyle\int\limits_1^3\left(2-2^x\right)\,\mathrm{\,d}x$}
	{$\displaystyle\int\limits_1^32^x\,\mathrm{\,d}x$}
	\loigiai{
		Hình phẳng gạch sọc trong hình vẽ được giới hạn bởi các đường $y=2^x,y=2,x=1$ và $x=3$ .\\
		Do đó diện tích hình phẳng gạch sọc trong hình vẽ bằng $\displaystyle\int\limits_1^3\left|2^x-2\right|\,\mathrm{\,d}x=\displaystyle\int\limits_1^3\left(2^x-2\right)\,\mathrm{\,d}x$.}
\end{ex}

\begin{ex}%Câu 11
	Cho tứ diện $ABCD.$ Gọi $M$ và $P$ lần lượt là trung điểm của các cạnh $AB$ và $CD.$ Đặt $\overrightarrow{BA}=\overrightarrow{b}$,  $\overrightarrow{AC}=\overrightarrow{c}$,  $\overrightarrow{AD}=\overrightarrow{d}$. Khẳng định nào sau đây là đúng?
	\choice
	{\True $\overrightarrow{MP}=\dfrac{1}{2}\left(\overrightarrow{c}+\overrightarrow{d}+\overrightarrow{b}\right)$}
	{$\overrightarrow{MP}=\dfrac{1}{2}\left(\overrightarrow{d}+\overrightarrow{b}-\overrightarrow{c}\right)$}
	{$\overrightarrow{MP}=\dfrac{1}{2}\left(\overrightarrow{c}+\overrightarrow{b}-\overrightarrow{d}\right)$}
	{$\overrightarrow{MP}=\dfrac{1}{2}\left(\overrightarrow{c}+\overrightarrow{d}-\overrightarrow{b}\right)$}
	\loigiai{
		\immini{
			\begin{align*}
				\overrightarrow{MP} & = \dfrac{1}{2}\left(\overrightarrow{MC}+\overrightarrow{MD}\right)                                        \\
				                    & =\dfrac{1}{2}\left(\overrightarrow{AC}-\overrightarrow{AM}+\overrightarrow{AD}-\overrightarrow{AM}\right) \\
				                    & =\dfrac{1}{2}\left(\overrightarrow{AC}+\overrightarrow{AD}-2\overrightarrow{AM}\right)                    \\
				                    & =\dfrac{1}{2}\left(\overrightarrow{AC}+\overrightarrow{AD}-\overrightarrow{AB}\right)                     \\
				                    & =\dfrac{1}{2}\left(\overrightarrow{c}+\overrightarrow{d}+\overrightarrow{b}\right).
			\end{align*}
		}
		{\begin{tikzpicture}[line join = round, line cap = round, thick, font = \small, scale = .7]
				\path
				(0:0) coordinate (B)
				+(0:5) coordinate (C)
				+(-70:3) coordinate (D)
				+(70:4) coordinate (A)
				($(A)!.5!(B)$) coordinate (M)
				($(C)!.5!(D)$) coordinate (P);
				\draw[dashed]
				(B)--(C) (M)--(P)
				;
				\draw
				(A)--(B)--(D)--(C)--cycle
				(A)--(D)
				;
				\foreach \x/\g in {B/-135,C/-45,D/-45,A/135,M/135,P/-45}
				\fill (\x) circle (1.5pt)
				+(\g:3mm) node {$\x$};
			\end{tikzpicture}}
	}
\end{ex}

\begin{ex}%Câu 12
	Thống kê điểm trung bình môn Toán của một số học sinh lớp $12$ được mẫu số liệu sau:\\
	\centerline{\begin{tabular}{|c|c|c|c|c|c|c|c|}
			\hline
			Khoảng điểm & $\left[6,5;7\right)$ & $\left[7;7,5\right)$ & $\left[7,5;8\right)$ & $\left[8;8,5\right)$ & $\left[8,5;9\right)$ & $\left[9;9,5\right)$ & $\left[9,5;10\right)$ \\
			\hline
			Tần số      & $8$                  & $10$                 & $16$                 & $24$                 & $13$                 & $7$                  & $4$                   \\
			\hline
		\end{tabular}}\\
	Phương sai của mẫu số liệu về điểm trung bình môn Toán của các học sinh đó là
	\choice
	{$0{,}616$}
	{$0{,}785$}
	{$0{,}78$}
	{\True $0{,}609$}
	\loigiai{
		Cỡ mẫu $n=8+10+16+24+13+7+4=82$ .\\
		Số trung bình của mẫu số liệu ghép nhóm là\\
		$$\bar{x}=\dfrac{8.6,75+10.7,25+16.7,75+24.8,25+13.8,75+7.9,25+4.9,75}{82}=\dfrac{333}{41}$$
		Phương sai của mẫu số liệu ghép nhóm là:\\
		$$s^2=\dfrac{1}{82}\left(8.6,75^2+10.7,25^2+16.7,75^2+24.8,25^2+13.8,75^2+7.9,25^2+4.9,75^2\right)-\left(\dfrac{333}{41}\right)^2 \approx 0,609$$.}
\end{ex}

\Closesolutionfile{ans}
\TNTF
\Opensolutionfile{ans}[ans/ans-DE-PNL-02-TF]

\begin{ex}%Câu 13
	Vẽ bảng biến thiên hàm số $f(x)=4\sin x\cos x+2x$ trên $\left[-\pi ;\pi\right]$.
	\choiceTF
	{Đạo hàm của hàm số đã cho là $f'(x)=4\sin 2x+2$}
	{\True Hàm số $y=f(x)$ có $4$ điểm cực trị thuộc $\left[-\pi ;\pi\right]$}
	{Hàm số $y=f(x)$ nghịch biến trên khoảng $\left(-2;-1\right)$}
	{\True Giá trị lớn nhất của $f(x)$ trên đoạn $\left[0;\dfrac{\pi}{2}\right]$ là $\dfrac{2\pi}{3}+\sqrt{3}$}
	\loigiai{
		\begin{itemchoice}
			\itemch $f(x)=4\sin x\cos x+2x=2\sin 2x+2x$ $\Rightarrow{f}'(x)=4\cos 2x+2$
			\itemch $f'(x)=0\Leftrightarrow\cos 2x=-\dfrac{1}{2}
				\Leftrightarrow\hoac{ & 2x=\dfrac{2\pi}{3}+k2\pi\\ & 2x=-\dfrac{2\pi}{3}+k2\pi}
				\Leftrightarrow \hoac{ & x=\dfrac{\pi}{3}+k\pi\\ & x=-\dfrac{\pi}{3}+k\pi},k\in\mathbb{Z}.$\\
			Trên đoạn $\left[-\pi ;\pi\right],f'(x)=0\Leftrightarrow x\in\left\{\dfrac{\pi}{3};-\dfrac{\pi}{3};\dfrac{2\pi}{3};-\dfrac{2\pi}{3}\right\}$. Vẽ bảng biến thiên của hàm số $y=f(x)$ trên đoạn $\left[-\pi ;\pi\right]$.
			\begin{center}
				\begin{tikzpicture}[>=stealth]
					\tkzTabInit[espcl=2.5,lgt=1.5,nocadre=false]
					{$x$/0.7,$f'(x)$/0.7,$f(x)$/2.1}
					{$-\pi$,$-\dfrac{2\pi}{3}$,$-\dfrac{\pi}{3}$,$\dfrac{\pi}{3}$,$\dfrac{2\pi}{3}$,$\pi$}
					\tkzTabLine{,+,0,-,0,+,0,-,0,+,}
					\tkzTabVar{-/ , +/ , -/ , +/ , -/ , +/ }
				\end{tikzpicture}
			\end{center}
			Từ bảng biến thiên ta thấy trên đoạn $\left[-\pi ;\pi\right]$ hàm số có 4 điểm cực trị.
			\itemch Nhận thấy $-\dfrac{2\pi}{3}< -2 < -\dfrac{\pi}{3} < -1$ nên từ bảng biến thiên trên ta có hàm số $y=f(x)$ không nghịch biến trên khoảng $\left(-2;-1\right)$.
			\itemch Dựa vào ta có $\max \limits _ {\left[0;\frac{\pi}{2}\right]} = f\left(\dfrac{\pi}{3}\right) = \dfrac{2\pi}{3}+\sqrt{3}$.
		\end{itemchoice}
	}
\end{ex}

\begin{ex}%Câu 14
	Một người điều khiển ô tô đang ở đường dẫn muốn nhập làn vào đường cao tốc. Khi ô tô cách điểm nhập làn $200$ m thì tốc độ của ô tô là $36$ (km/h). Hai giây sau đó, ô tô bắt đầu tăng tốc với tốc độ $v(t)=at+b$ ($a,b\in\mathbb{R},a>0$), trong đó $t$ là thời gian tính bằng giây kể từ khi bắt đầu tăng tốc. Biết rằng ô tô nhập làn cao tốc sau $12$ giây và duy trì sự tăng tốc trong $24$ giây kể từ khi bắt đầu tăng tốc. Sau $24$ giây đó ô tô duy trì tốc độ cao nhất trong thời gian còn lại trên cao tốc.\\
	\centerline{\includegraphics[height=5cm]{images/2.14}}
	\choiceTF
	{\True Quãng đường ô tô đi được từ khi bắt đầu tăng tốc đến khi nhập làn là $180$ m}
	{Vận tốc của ô tô tại thời điểm nhập làn là $72$ (km/h)}
	{Quãng đường mà ô tô đi được trong thời gian $30$ giây kể từ khi ô tô cách điểm nhập làn $200$ m là $620$ m}
	{Sau 24 giây kể từ khi tăng tốc, ô tô duy trì tốc độ cao nhất trong vòng $5$ giây thì phát hiện chướng ngoại vật cách đó $300$ m. Người điều khiển lập tức đạp phanh và ô tô chuyển động chậm dần đều với $a(t)=-3$ (m/s$^2$). Khi đó ô tô dừng lại cách chứng ngại vật $10$ m}
	\loigiai{
		\begin{itemchoice}
			\itemch Ta có $36$ (km/h)$=10$ (m/s).\\
			Vận tốc ô tô lúc đầu là $10$ (m/s) nên trong $2$ giây ô tô đi được quãng đường $2 \cdot 10=20$ m\\
			Quãng đường ô tô đi được từ khi bắt đầu tăng tốc đến khi nhập làn là $200-20=180$ m.
			\itemch Ôtô nhập làn sau $12$ giây tăng tốc nên
			$$s(12)=180 \Leftrightarrow \displaystyle\int\limits_0^{12}{v(t)dt}=180 \Leftrightarrow \dfrac{1}{2}.a.12^2+12b=180 \Leftrightarrow 72a+12b=180.$$
			Mặt khác: $v(0)=10$ nên $b=10$ suy ra $a=\dfrac{5}{6}$\\
			Do đó vận tốc ô tô tại thời điểm nhập làn là $v(12)=20$ (m/s) $=72$ (km/h).
			\itemch Quãng đường ô tô đi được sau $24$ giây tăng tốc là $s(24)=\displaystyle\int\limits_0^{24}{v(t)dt}=\displaystyle\int\limits_0^{24}{\left(at+b\right)dt}=480$ m\\
			Vận tốc tại thời điểm 24 giây từ lúc tăng tốc là $v\left(24\right)=\dfrac{5}{6}.24+10=30$ (m/s)\\
			Quãng đường $4$ giây kể từ lúc vận tốc đạt lớn nhất là $30.4=120$ m\\
			Vậy tổng quãng đường đi được sau $30$ giây kể từ khi ô tô cách điểm nhập làn $200$ m là\\
			$20+480+120=620$ m.
			\itemch Vận tốc tốc tại thời điểm vượt chướng ngại vật là $30$ (m/s)\\
			Vận tốc từ khi bắt đầu phanh là $v=\displaystyle\int a dt=\displaystyle\int -3dt=-3t+C$\\
			Vì vận tốc khi phanh là $30$ (m/s) nên $C=30$ do đó vận tốc là $v=-3t+30$\\
			Khi xe dừng thì vận tốc bằng $0$ nên ta có $-3t+30=0\Leftrightarrow t=10$ (giây)\\
			Quãng đường đi được đến khi xe dừng là $S=\displaystyle\int\limits_0^{10}{\left(-3t+30\right)dt}=150$ m\\
			Vậy ôtô cách chướng ngại vật $300-150=150$ m.
		\end{itemchoice}
	}
\end{ex}

\begin{ex}%Câu 15
	Một loại sản phẩm do hai nhà máy số I, số II cùng sản xuất. Tỷ lệ phế phẩm của các nhà máy I, II lần lượt là $0,04$ ; $0,03$ . Trong một lô sản phẩm để lẫn lộn $80$ sản phẩm của nhà máy số I và $120$ sản phẩm nhà máy số II. Một khách hàng lấy ngẫu nhiên 1 sản phẩm từ lô hàng đó.
	\choiceTF
	{\True Số phần tử của không gian mẫu là $200$}
	{\True Xác suất để lấy được sản phẩm tốt là $\dfrac{483}{500}$}
	{Biết sản phẩm lấy được là phế phẩm, xác suất sản phẩm được sản xuất từ nhà máy I là $\dfrac{8}{19}$}
	{Khả năng lấy được sản phẩm không tốt của nhà máy II là thấp hơn nhà máy I}
	\loigiai{
		\begin{itemchoice}
			\itemch Số phần tử của không gian mẫu là $n\left(\Omega\right)=80+120=200$.
			\itemch Gọi $B$ là biến cố lấy được sản phẩn không tốt.\\
			Xác suất để lấy được sản phẩm tốt là $P(\bar{B})=\dfrac{80}{200}.0,96+\dfrac{120}{200}.0,97=\dfrac{483}{500}$.
			\itemch Gọi $A$ là biến cố lấy được sản phẩm của nhà máy I.\\
			Khi đó $P\left(A|B\right)=\dfrac{P\left(A\cap B\right)}{P(B)}=\dfrac{P\left(B|A\right)P(A)}{P(B)}=\dfrac{0,04 \cdot \dfrac{80}{200}}{1-\dfrac{483}{500}}=\dfrac{8}{17}$.
			\itemch Xác suất để lấy được sản phẩm không tốt ở máy I là $P(AB)=P(B|A) \cdot P(A)=0,04 \cdot \dfrac{80}{200}=\dfrac{8}{100}=0,016$\\
			Xác suất để lấy được sản phẩm không tốt ở máy II là $P\left(\overline{A}B\right)=P(B|\overline{A}) \cdot P(\overline{A})=0,03 \cdot \dfrac{120}{200}=\dfrac{3}{100}=0,018$\\
			Vậy khả năng lấy được sản phẩm không tốt ở máy II cao hơn máy I.
		\end{itemchoice}
	}
\end{ex}

\begin{ex}%Câu 16
	Các thiên thạch có đường kính lớn hơn $140m$ và có thể lại gần Trái Đất ở khoảng cách nhỏ hơn $7500000km$ được coi là những vật thể có khả năng va chạm gây nguy hiểm cho Trái Đất. Để theo dõi những thiên thạch này, các nhà nghiên cứu của trung tâm Vũ Trụ Nasa đã thiết lập các trạm quan sát các vật thể bay gần Trái Đất. Giả sử có một hệ thống quan sát có khả năng theo dõi các vật thể ở độ cao không vượt quá $4600km$ so với mực nước biển. Coi Trái Đất là khối cầu có bán kính $6400km$ . Chọn hệ trục tọa độ $Oxyz$ trong không gian có gốc $O$ tại tâm Trái Đất và đơn vị độ dài trên mỗi trục tọa độ là $1000$ km. Một thiên thạch (coi như một hạt) chuyển động với tốc độ $v_1=2\sqrt{2}.10^3$ (km/h) không đổi theo đường thẳng xuất phát từ điểm $M\left(0;5;12\right)$ đến $N\left(12;5;0\right)$\\
	\centerline{\includegraphics[width=.4\textwidth]{images/2.16.png}}
	\choiceTF
	{\True Khoảng cách thiên thạch gần với trái đất nhất có độ dài bằng $3449km$ (Kết quả làm tròn đến hàng đơn vị)}
	{\True Các nhà nghiên cứu của trung tâm vũ trụ Nasa đưa ra giả thiết nếu lúc thiên thạch đang ở vị trí $M$ bất ngờ đổi hướng và lao xuống trái đất với phương thẳng thì quãng đường dài nhất nó có thể va chạm với trái đất là $11315$ km (Kết quả làm tròn đến hàng đơn vị)}
	{Tại thời điểm thiên thạch đang ở vị trí $M$ thì có 2 vệ tinh đang ở vị trí $A\left(-6;\,-5;\,-6\right)$ và $B\left(7;\,-6;\,7\right)$ có vận tốc khác nhau di chuyển trong mặt phẳng trung trực của $MN$ và luôn cách trái đất với khoảng cố định. Khoảng cách xa nhất của 2 vệ tinh có thể đạt là $18412km$ (Kết quả làm tròn đến hàng đơn vị)}
	{\True Nếu vệ tinh $A$ đi với vận tốc $v_2=\dfrac{\pi\sqrt{97}}{3}.10^3$ (km/h) thì sẽ va chạm với thiên thạch}
	\loigiai{
		\begin{itemchoice}
			\itemch Phương trình đường thẳng $MN$ là: $\heva{& x=12t\\&y=5\\&z=12-12t}$\\
			Khoảng cách thiên thạch gần tâm Trái Đất nhất là $d(O,MN)=\dfrac{|[\vec{OM},\vec{MN}]|}{|\vec{MN}|}=\sqrt{97}$\\
			Suy ra khoảng cách thiên thạch gần Trái Đất nhất là $d(O,MN)-R=\sqrt{97}.1000-6400\approx 3449$ (km)\\
			\itemch Quãng đường dài nhất thiên thạch va chạm trái đất là $MK=\sqrt{OM^2-R^2} \approx \sqrt{13^2-6,4^2}.10^3\approx 11315$ (km)\\
			\itemch Phương trình mặt phẳng trung trực $MN$ là $x-z=0$\\
			Hai vệ tinh $A$, $B$ và tâm $O$ đều nằm trên mặt phẳng trung trực của $MN$ nên khoảng cách xa nhất của hai vệ tinh là $AB=OA+OB=\sqrt{6^2+5^2+6^2}+\sqrt{7^2+6^2+7^2}=\sqrt{97}+\sqrt{134}\approx 21425$ (km)\\
			\itemch Giả sử vệ tinh $A$ có thể va chạm với thiên thạch. Gọi vị trí va chạm là $H$. Khi đó $H$ là giao điểm của $MN$ với mặt phẳng trung trực của nó (vì $A$ nằm trên đó), do đó $H$ là trung điểm $MN$ và có toạ độ $H\left(6;5;6\right)$\\
			Ta có $AH=\sqrt{12^2+10^2+12^2}=\sqrt{388}$ (nghìn km)\\
			Khi đó: $\cos\widehat{AOH}=\dfrac{2OA^2-AH^2}{2OA \cdot OH}=\dfrac{2 \cdot 97-388}{2\sqrt{97} \cdot \sqrt{97}}=-1\Rightarrow\widehat{AOH}=180^\circ$\\
			Độ dài cung trên $AH$ là $l_{AH}=\pi\sqrt{97}.10^3$ (km)\\
			Thời gian vệ tinh di chuyển từ $A$ đến $H$ là $t_2=\dfrac{l_{AH}}{v_2}=\dfrac{\left(\pi \cdot \sqrt{97}.10^3\right)}{\left(\pi \cdot \dfrac{\sqrt{97}}{3}.10^3\right)}=3$ (giờ)\\
			Thời gian thiên thạch đi từ $M$ tới $H$ là $t_{MH}=\dfrac{MH}{v_1}=\dfrac{\sqrt{72}.10^3}{2\sqrt{2}.10^3}=3$ (giờ)\\
			Vậy vì $t_1=t_2$ nên thiên thạch và vệ tinh $A$ sẽ va chạm với nhau.
		\end{itemchoice}
	}
\end{ex}
\Closesolutionfile{ans}
\TNSA
\Opensolutionfile{ans}[ans/ans-DE-PNL-02-SA]
\begin{ex}%Câu 17
	Cho hình lăng trụ đứng $ABC.\,A'{B}'{C}'$ có đáy $ABC$ là tam giác đều cạnh bằng$\sqrt{2}$ và độ dài cạnh$B{A}'=\sqrt{6}$ . Hãy tính khoảng cách giữa hai đường thẳng $A'B$ và $B'C$ (Kết quả làm tròn đến hàng phần trăm)?\\
	\centerline{\begin{tikzpicture}[scale=1, font=\footnotesize, line join=round, line cap=round, >=stealth,thick]
			\def\ac{4} % cạnh AC
			\def\ab{2} % cạnh AB
			\def\h{4} % chiều cao
			\def\gocA{50} % góc A của đáy
			\path
			(0,0) coordinate (A)
			(\ac,0) coordinate (C)
			(-\gocA:\ab) coordinate (B)
			($(A)+(90:\h)$) coordinate (A')
			($(B)-(A)+(A')$) coordinate (B')
			($(C)-(A)+(A')$) coordinate (C');
			\draw (B)--(A')--(A)--(B)--(C)--(C')--(A')--(B')--(C') (B)--(B')--(C);
			\draw[dashed] (A)--(C);
			\foreach \x/\g in {A/180,B/-90,C/0,A'/180,C'/0,B'/-140}\fill[red] (\x) circle (1pt)+(\g:3mm) node[black]{$ \x $};
		\end{tikzpicture}}
	\shortans{0,67}
	\loigiai{
	Gọi $M,\,N$ lần lượt là trung điểm của $AC$ và $A'B$ suy ra $MN\parallel{B}'C$ nên $B'C\parallel\left(A'BM\right)$ .\\
	Khi đó: $d\left(\,B'C,\,A'B\right)=d\left(\,B'C,\,\left(A'BM\right)\right)=d\left(C,\,\left(A'BM\right)\right)=d\left(A,\,\left(A'BM\right)\right)$ .\\
	Kẻ $AH\bot{A}'M$ tại $H$ thì ta có $\left\{\begin{aligned}
			 & BM\perp AC  \\
			 & BM\perp AA' \\
		\end{aligned}\right.\Rightarrow BM\perp\left(AC{C}'{A}'\right)$ .\\
	Khi đó $\left\{\begin{aligned}
			 & AH\bot{A}'M \\
			 & AH\perp BM  \\
		\end{aligned}\right.\Rightarrow AH\perp\left(A'BM\right)$ nên $d\left(A,\,\left(A'BM\right)\right)=AH$ .\\
	Ta có $AM=\dfrac{1}{2}AC=\dfrac{\sqrt{2}}{2}$ ; $A{A}'=\sqrt{A'{B^2}-A{B^2}}=2$ .\\
	Trong $\Delta A{A}'K$ ta có: $\dfrac{1}{A{H^2}}=\dfrac{1}{A{A'^2}}+\dfrac{1}{A{M^2}}\Rightarrow d\left(\,B'C,\,A'B\right)=AH=\dfrac{A{A}'.\,AM}{\sqrt{A{A'^2}+A{M^2}}}=\dfrac{2}{3}\approx 0,67$
	}
\end{ex}

\begin{ex}%Câu 18
	Có năm địa điểm $A,B,C,D,E$. Một số địa điểm có đường đi tới nhau môt tả bằng các cạnh với độ dài quãng đường tính theo kilomet cho bởi số gắn với cạnh đó như hình vẽ. Một người đưa thư xuất phát từ bưu điện ở vị trí $C$ , cần đi qua tất cả các đường (mỗi đường đi qua ít nhất một lần) và sau đó phải trở về vị trí ban đầu $C$. Tổng số kilomet mà người đưa thư phải đi nhỏ nhất là bao nhiêu?\\
	\centerline{
		\begin{tikzpicture}[declare function={r=3;},thick]
			\path
			(0,0) coordinate (B)
			(0:r) coordinate (C)
			(-90:r) coordinate (D)
			($(C)+(D)-(B)$) coordinate (E)
			($(B)!.5!(C) +(90:2)$) coordinate (A)
			;
			\draw (B) rectangle (E) (E)--(B)node[midway,below]{$5$}--(A) node[midway,left]{$1$}--(C)node[midway,above]{$2$}
			;
			\path (B)--(D) node[left,midway]{$3$}
			(B)--(C) node[left,midway]{$3$}
			(E)--(D) node[below,midway]{$6$}
			(E)--(C) node[right,midway]{$10$};
			\foreach \x/\g in {A/90,B/180,C/0,D/-90,E/-90}\draw[fill=white] (\x) circle (1pt)+(\g:3mm) node{$\x$};
		\end{tikzpicture}
	}
	\shortans{38}
	\loigiai{
		Đi từ $C$ đến $E$ theo đường đi Euler $CABCEBDE$ dài: $2+1+3+3+5+6+10=30$ (km)\\
		Đi từ $E$ đến $C$ với quảng đường ngắn nhất dài: $5+3=8$ (km)\\
		Vậy tổng số kilomet mà người đưa thư phải đi nhỏ nhất là $30+8=38$ (km).}
\end{ex}

\begin{ex}%Câu 19
	Hệ thống định vị toàn cầu GPS (Global Positioning System) là một hệ thống cho phép xác định vị trí của một vật thể trong không gian. Trong cùng một thời điểm, vị trí của một điểm $M$ trong không gian sẽ được xác định bởi bốn vệ tinh cho trước nhờ các bộ thu phát tín hiệu đặt trên các vệ tinh. Giả sử trong không gian với hệ tọa độ $Oxyz$ , tỉ lệ độ dài trên các trục là $10$ km tính cho một đơn vị tỉ lệ trên mỗi trục có 4 vệ tinh lần lượt đặt tại các điểm $A(-1;2;-1)$ ,$B(1;4;0)$ ,$C(3;0;9)$ và $D(7;10;-1)$ . Ở một thời điểm cả bốn vệ tinh bắn tín hiệu về điểm $M$ và đo được độ dài $MA=6,\,MB=3,\,MC=10,\,MD=6$ . Ngay sau đó $10$ giây, cả bốn vệ tinh lại bắn tin hiệu về vật $M$ và đo được $MA=3,\,MB=4,\,MC=\sqrt{85},\,MD=\sqrt{137}$. Nếu coi như vật $M$ chuyển động thẳng đều thì tốc độ của vật bằng bao nhiêu (đơn vị: m/s và kết quả làm tròn đến hang đơn vị)? (Bỏ qua khoảng thời gian phát và thu tín hiệu)
	\shortans{6403}
	\loigiai{
		Gọi $M\left(a;\,b;\,c\right)$ .\\
		Tại thời điểm đầu bốn vị tinh bắn tín hiệu về điểm $M$ . Khi đó ta có hệ phương trình sau:\\
		$\Leftrightarrow\left\{\begin{matrix}
				{\left(a+1\right)^2}+\left(b-2\right)^2+\left(c+1\right)^2=36\,\,\,\,\,\,\,\,\,\,\,\,\,\,(1)               \\
				{\left(a-1\right)^2}+\left(b-4\right)^2+c^2=9\,\,\,\,\,\,\,\,\,\,\,\,\,\,\,\,\,\,\,\,\,\,\,\,\,\,\,\,\,(2) \\
				{\left(a-3\right)^2}+b^2+\left(c-9\right)^2=100\,\,\,\,\,\,\,\,\,\,\,\,\,\,\,\,\,\,\,\,\,\,\,(3)           \\
				{\left(a-7\right)^2}+\left(b-10\right)^2+\left(c+1\right)^2=36\,\,\,\,\,\,\,\,\,\,\,(4)                    \\
			\end{matrix}\right.\Leftrightarrow\left\{\begin{aligned}
				 & (2)-(1):-4a-4b-2c=-38  \\
				 & (3)-(1):-8a+4b-20c=-20 \\
				 & (4)-(1):-14a-14b=-144  \\
			\end{aligned}\right.$\\
		Giải hệ phương trình này, ta tìm được $a~=~3,b=6,c=1\Rightarrow M\left(3;\,6;\,1\right)$\\
		Sau đó 10 giây, về tinh lại bán tín hiệu về điểm $M$ . Gọi $M'\left(x;\,y;\,z\right)$ khi đó ta có:\\
		$\Leftrightarrow\left\{\begin{matrix}
				{\left(x+1\right)^2}+\left(y-2\right)^2+\left(z+1\right)^2=9\,\,\,\,\,\,\,\,\,\,\,\,\,\,(1)     \\
				{\left(x-1\right)^2}+\left(y-4\right)^2+z^2=16\,\,\,\,\,\,\,\,\,\,\,\,\,\,\,\,\,\,\,\,\,\,\,(2) \\
				{\left(x-3\right)^2}+y^2+\left(z-9\right)^2=85\,\,\,\,\,\,\,\,\,\,\,\,\,\,\,\,\,\,\,\,\,\,(3)   \\
				{\left(x-7\right)^2}+\left(y-10\right)^2+\left(z+1\right)^2=137\,\,\,\,(4)                      \\
			\end{matrix}\right.\Leftrightarrow\left\{\begin{aligned}
				 & (2)-(1):-4x-4y-2z=-4  \\
				 & (3)-(1):-8x+4y-20z=-8 \\
				 & (4)-(1):-14x-14y=-16  \\
			\end{aligned}\right.$\\
		Giải hệ phương trình này, ta tìm được $x~=~1,y=0,z=0\Rightarrow{M}'\left(1;\,0;\,0\right)$\\
		Khi đó khoảng cách $M{M}'=\sqrt{\left(1-3\right)^2+\left(0-6\right)^2+\left(0-1\right)^2}=\sqrt{41}$\\
		Do vật chuyển động thẳng đều nên vận tốc của vật là: $v_M=\dfrac{M{M}'}{t}=\dfrac{\sqrt{41}{10^4}}{10}=6403$ (m/s)
	}
\end{ex}
\begin{ex}%Câu 20
	Một ly trà sữa dạng hình nón cụt, có đường kính đáy ly 6 cm, đường kính miệng ly 9 cm, chiều cao 13,4 cm, ở miệng ly có sử dụng một nắp đậy có hình dạng nửa mặt cầu và ở đỉnh của nửa mặt cầu này có một hình tròn có đường kính 2 cm để cắm ống hút, mặt phẳng chứa hình tròn này song song với mặt phẳng chứa miệng ly (tham khảo hình vẽ sau).
	\begin{center}
		\includegraphics[width=.3\textwidth]{images/2.20a.png} \includegraphics[width=.22\textwidth]{images/2.20b.png}
		\includegraphics[width=.3\textwidth]{images/2.20c.png}
	\end{center}
	Chọn hệ trục $Oxy$ (đơn vị trên trục là centimet) với trục $Ox$ đi qua tâm của 2 đáy hình nón cụt và gốc tọa độ $O$ trùng với tâm của đáy lớn như hình vẽ trên. Tính thể tích bên trong của ly bao gồm cả thể tích của nắp (Kết quả làm tròn kết quả đến hàng đơn vị).
	\shortans{790}
	\loigiai{
	Thể tích bên trong của ly không bao gồm nắp là: $V_1=\pi{\displaystyle\int\limits_{-13,4}^0\left(\dfrac{1,5x+60,3}{13,4}\right)^2}dx\approx 600$ cm3\\
	Thể tích bên trong của ly bao gồm cả thể tích của nắp là: $V=V_1+\left(V_2-V_3\right)$ trong đó $V_2$ là nửa thể tích của khối cầu và $V_3$ là thể tích của chỏm cầu.\\
	Nửa thể tích khối cầu là: $V_2=\dfrac{1}{2}.\dfrac{4}{3}.\pi .R^3=\dfrac{243\pi}{4}$\\
	Thể tích chỏm cầu: $V_3=\pi\displaystyle\int\limits_{R-h}^R{\left(\sqrt{R^2-x^2}\right)^2dx}$ $\Leftrightarrow{V_3}=\pi\displaystyle\int\limits_{R-h}^R{\left(R^2-x^2\right)dx\Leftrightarrow{V_3}=\,}\pi\left.\left(R^2x-\dfrac{x^3}{3}\right)\right|_{R-h}^R$\\
	$\Leftrightarrow{V_3}=\pi\left[\left(R^3-\dfrac{R^3}{3}\right)-\left(R^2\left(R-h\right)-\dfrac{\left(R-h\right)^3}{3}\right)\right]$ $\Leftrightarrow{V_3}=\pi{h^2}\left(R-\dfrac{h}{3}\right)$ .\\
	Thay số ta có: $h=4,5-\dfrac{\sqrt{77}}{2}$ ; $R=4,5$ $\Rightarrow{V_3}=\dfrac{\left(729-83\sqrt{77}\right)\pi}{12}$\\
	Thể tích bên trong của ly bao gồm cả thể tích của nắp là:\\
	$V=600+\left[\dfrac{243\pi}{4}-\dfrac{\left(729-83\sqrt{77}\right)\pi}{12}\right]\approx 790$ ml.}
\end{ex}

\begin{ex}%Câu 21
	Một doanh nghiệp dự định sản xuất không quá $400$ sản phẩm. Nếu doanh nghiệp sản xuất $x$ sản phẩm $\left(1\le x\le 400\right)$ thì doanh thu nhận được khi bán hết số sản phẩm đó được biểu diễn bởi công thức là $F(x)=x^3-1999x^2+1001000x+250000$ (đồng). Trong đó chi phí vận hành máy móc cho mỗi sản phẩm là $G(x)=\dfrac{100000x}{\dfrac{3}{2}x+1}$ (đồng). Tổng chi phí mua nguyên vật liệu được biểu diễn bởi hàm $H(x)=2x^3+100000x-50000$ (đồng) nhưng do doanh nghiệp đó mua nguyên vật liệu với số lượng lớn nên được giảm $1\%$ cho 200 sản phẩm đầu tiên doanh nghiệp sản xuất và giảm $2\%$ cho sản phẩm tiếp theo. Doanh nghiệp cần sản xuất bao nhiêu sản phẩm để lợi nhuận thu được là lớn nhất?
	\shortans{184}
	\loigiai{
	Lợi nhuận $P(x)$ được tính bằng doanh thu trừ đi tổng chi phí: $P(x)=F(x)-xG(x)-H(x)$ .\\
	Khi $x\le 200$ thì chi phí mua nguyên liệu là:
	$$H_1(x)=0,99H(x)=0,99\left(2x^3+100000x-50000\right) \text{ (đồng).}$$
	Khi $x>200$ thì chi phí mua nguyên liệu là:
	$$H_2(x)=0,99H(200)+0,98[H(x)-H(200)]=0,01H(200)+0,98H(x) \text{ (đồng).}$$
	Xét đồng thời 2 trường hợp:\\
	Trường hợp 1: Với $1\le x\le 200$ thì ta có lợi nhuận thu được là:
	\begin{align*}
		P_1(x) & =F(x)-xG(x)-H_1(x)                                                 \\
		       & =-0,98x^3-1999x^2+902000x+299500-\dfrac{100000x^2}{\frac{3}{2}x+1} \\
	\end{align*}
	Ta có: $P_1'(x)=-2,94{x^2}-3998x-\dfrac{600000x^2+800000x}{\left(3x+2\right)^2}+902000$\\
	Phương trình $P_1'(x)=0$ có nghiệm $x=184,03\in\left(1;200\right)$ .\\
	Ta thấy $\max\limits_{\left[1;200\right]} P_1(x)=80037062,09$ tại $x=184,03$ .\\
	Trường hợp 2: Với $201\le x\le 400$ ta có lợi nhuận thu được là:
	\begin{align*}
		P_2(x) & =F(x)-xG(x)-H_2(x)                                                 \\
		       & =-0,96x^3-1999x^2+903000x+609500-\dfrac{100000x^2}{\frac{3}{2}x+1}
	\end{align*}
	Ta có $P_2'(x)=-2,88{x^2}-3998x+903000-\dfrac{600000x^2+800000x}{\left(3x+2\right)^2}$\\
	Phương trình $P_2'(x)=0$ không có nghiệm thuộc $\left(201;400\right)$ .\\
	Suy ra $\max\limits_{\left[201;400\right]} P_2(x)=P(201)\approx 7,959.10^8$\\
	Vậy doanh nghiệp cần sản xuất 184 sản phẩm thì lợi nhuận thu được là lớn nhất.
	}
\end{ex}

\begin{ex}%Câu 22
	Có 10 học sinh làm bài kiểm tra xác suất thống kê, trong đó có 2 học sinh giỏi (trả lời được 100\% các câu hỏi), 3 học sinh khá (trả lời được 80\% các câu hỏi), 5 học sinh trung bình (trả lời được 50\% các câu hỏi). Bài kiểm tra có 4 câu hỏi được lấy ngẫu nhiên từ 20 câu hỏi. Giáo viên chọn ngẫu nhiên một bài làm của học sinh để chấm điểm. Xác suất bài làm đó trả lời được cả 4 câu hỏi là bao nhiêu (kết quả làm tròn đến hàng phần trăm)?

	\shortans{0,33}
	\loigiai{
		Gọi $A$ là biến cố chọn được bài làm trả lời được cả 4 câu hỏi.\\
		Gọi $B_1,B_2,B_3$ lần lượt là biến cố chọn được bài làm của học sinh giỏi, khá, trung bình.\\
		Khi đó: $P\left(B_1\right)=\dfrac{2}{10},P\left(B_2\right)=\dfrac{3}{10},P\left(B_3\right)=\dfrac{5}{10}$ .\\
		Học sinh giỏi trả lời được 100\% các câu hỏi, nghĩa là trả lời được 20 câu suy ra $P\left(A|B_1\right)=\dfrac{C_{20}^4}{C_{20}^4}.$\\
		Học sinh khá trả lời được 80\% các câu hỏi, nghĩa là trả lời được 16 câu suy ra $P\left(A|B_2\right)=\dfrac{C_{16}^4}{C_{20}^4}.$\\
		Học sinh trung bình trả lời được 50\% các câu hỏi, nghĩa là trả lời được 10 câu suy ra\\
		$P\left(A|B_3\right)=\dfrac{C_{10}^4}{C_{20}^4}$ .\\
		Theo công thức xác suất toàn phần ta có:\\
		$P(A)=P\left(B_1\right).P\left(A|B_1\right)+P\left(B_2\right).P\left(A|B_2\right)+P\left(B_3\right).P\left(A|B_3\right)$\\
		$=\dfrac{2}{10}.\dfrac{C_{20}^4}{C_{20}^4}+\dfrac{3}{10}.\dfrac{C_{16}^4}{C_{20}^4}+\dfrac{5}{10}.\dfrac{C_{10}^4}{C_{20}^4}=\dfrac{108}{323}\approx 0,33$}
\end{ex}
\Closesolutionfile{ans}
\Closesolutionfile{ansbook}
\inputansbox{6,2,3}{ans/ans-DE-PNL-02-T,ans/ans-DE-PNL-02-TF,ans/ans-DE-PNL-02-SA}